\section{Strict exterior zeros obstruct the actual full character matrix}
\label{sec:character-exterior-obstruction}

The sufficient half-plane premise in MA-MO-022 is not a
disposable device for its inverse estimate. A strictly exterior
zero obstructs boundedness of the original full common-sum
action itself, even after any nonzero compact smoothing.
The proof uses genuine translated smooth inputs and an actual
bounded output functional. It does not replace the original
inverse by a narrow multiplier continuation.

Fix a number field \(K\), a nontrivial primitive finite-order
Hecke character \(\eta\), the structural place set and actual
source conventions of (MA-MO.470), and
\[
 \begin{gathered}
 H=H_{\Sigma,a,s},\qquad 0<a<1/2,\qquad s\ge0,\\
 C_gf(x)=\int g(t)f(x-t)\,dt,\qquad
 U_tf(x)=f(x-t),\qquad Rf(x)=f(-x),\\
 \phi=1\ \hbox{sufficiently far left},\qquad
 \phi=0\ \hbox{sufficiently far right},\qquad
 \phi'\in C_c^\infty,\qquad q=1-\phi .
 \end{gathered}
 \tag{MA-MO.550}
\]
The cutoff is real and smooth; no monotonicity or symmetry
is assumed. In this section there is \emph{no} zero-free
half-plane hypothesis. Put
\[
 \begin{aligned}
 A_\omega(z)&=E_{\omega,S}(1/2+z)L_K(1/2+z,\omega),\\
 F_{\omega,y}(x)&=
   \sum_{\mathfrak a\ {\rm prime\ to}\ S}
     \mu_K(\mathfrak a)c_{\omega,\mathfrak a}
       (N\mathfrak a)^{-1/2}y(x+\log N\mathfrak a),
       \qquad \omega=\eta,\bar\eta .
 \end{aligned}
 \tag{MA-MO.551}
\]
These are the original wide Euler inverses of MA-CC.1.
Write \(\mathfrak S_\eta,\mathfrak S_{\bar\eta}\) for the
actual inverse Fourier pair on the two local source sectors.
The original smooth common-space action is
\[
 \begin{aligned}
 \mathscr D^{\,\rm raw}_{\eta,g}y
   &=C_g\{\vartheta_\eta H_{\eta,y}-y\},\\
 H_{\eta,y}
   &=qF_{\eta,y}
       +\mathfrak S_{\bar\eta}(qF_{\bar\eta,Ry}),
                         \qquad y\in\mathcal E .
 \end{aligned}
 \tag{MA-MO.552}
\]
It is \(C_g(P_+-P_-)\) on the actual nontrivial common
sector, where \(P_-=1\). All finite Fourier phases are
those of the original source.

\paragraph{MA-MO-024 (actual-source strict-exterior obstruction).}
If a completed nontrivial zero \(\rho\) of \(L_K(s,\eta)\)
satisfies
\[
                  |\Re\rho-1/2|>a,
 \tag{MA-MO.553}
\]
then for every nonzero \(g\in C_c^\infty(\mathbb R)\)
the map (MA-MO.552) has no bounded extension \(H\to H\).
The conclusion holds for every polynomial exponent \(s\ge0\).
For a right-exterior zero there are a fixed
\(h\in C_c^\infty(\mathbb R)\) and \(t_n\to+\infty\)
such that the genuine smooth inputs have
\[
       \frac{\|\mathscr D^{\,\rm raw}_{\eta,g}U_{t_n}h\|_H}
            {\|U_{t_n}h\|_H}\longrightarrow\infty .
 \tag{MA-MO.554}
\]
For a left-exterior zero the corresponding reflected
translation sequence gives the conclusion. Both inputs
and outputs in these sequences belong to the original
smooth common sector.

In particular, existence of a bounded extension for even
one nonzero test implies that the entire nontrivial divisor
lies in the \emph{closed} strip
\[
                   |\Re\rho-1/2|\le a .
 \tag{MA-MO.555}
\]
No conclusion about zeros on its boundary is asserted.
There is no exception when \(\mathcal Mg(\rho-1/2)=0\).
The theorem also applies to nonzero tests with
\(g(0)=\mathcal Mg(1/2)=\mathcal Mg(-1/2)=0\).

\begin{proof}
\emph{The source domain before any narrow estimate.}
For \(y\in\mathcal E\), the series (MA-MO.551) and
every derivative converge locally absolutely. Its right
tail and all derivatives decrease with every exponential
weight, by the wide Euler estimates at arbitrarily large
weights. Multiplication by \(q\) kills the other tail.
Thus \(qF_{\omega,y}\in\mathcal E_\omega\), without
moving any inverse contour into the critical strip.
The actual Schwartz-cup inclusion (MA-MO.476) puts it
in \(\mathcal T_\omega\). The two Fourier maps in
(MA-MO.552) consequently act on their original domains.
Nontrivial-sector Poisson summation puts
\(\vartheta_\eta H_{\eta,y}\) in \(I_\eta\subset\mathcal E\).
Therefore (MA-MO.552) is a genuine map
\(\mathcal E\to\mathcal E\) independently of the bounded
extension under investigation.

\emph{Support and elementary wide estimates.}
For now suppose
\(\lambda=\rho-1/2\) satisfies \(\Re\lambda>a\).
Use bilinear, complex-linear output functionals
\[
                  \ell_v(f)=\int f(x)v(x)\,dx,
                         \qquad v\in C_c^\infty(\mathbb R).
 \tag{MA-MO.556}
\]
Each is continuous on \(H\) by weighted Cauchy--Schwarz.
There is no complex conjugation of \(g\) in its transpose:
\[
 \ell_v(C_gf)=\ell_{v_g}(f),\qquad
 v_g=C_{g^\vee}v,\qquad g^\vee(t)=g(-t).
 \tag{MA-MO.557}
\]
The choices of \(h\) and \(v\) that detect \(\lambda\)
will be made below; first take any compact smooth \(h,v\).
Write
\[
 \begin{gathered}
 \operatorname{supp}h\subset[\alpha,\beta],\qquad
 \operatorname{supp}v_g\subset[c,d],\qquad
 \phi(x)=1\quad(x\le\ell),\\
 F=F_{\eta,h},\quad
 H_h(w)=\int h(x)e^{wx}\,dx,\quad
 V_g(w)=\int v_g(x)e^{wx}\,dx,\\
 W(y)=\sum_{\mathfrak a\ {\rm prime\ to}\ S}
       c_{\eta,\mathfrak a}(N\mathfrak a)^{-1/2}
                            v_g(y-\log N\mathfrak a).
 \end{gathered}
 \tag{MA-MO.558}
\]
Both sums defining \(F\) and \(W\) are locally finite
because the two tests are compactly supported and there
are only finitely many integral ideals of bounded norm.
In particular
\[
           F(x)=0\ (x>\beta),\qquad
           W(y)=0\ (y<c).
 \tag{MA-MO.559}
\]
For every fixed \(b>1/2\), absolute ideal Euler convergence
also gives
\[
             |F(x)|\le C_b e^{-bx},\qquad
             |W(y)|\le C_b' e^{by}.
 \tag{MA-MO.560}
\]
For example, factor \(e^{-bx}(N\mathfrak a)^{-b}\)
out of \(h(x+\log N\mathfrak a)\); the remaining sum
is bounded by \(\zeta_K(b+1/2)\). The other estimate
uses \(e^{-b(y-\log N\mathfrak a)}v_g(y-\log N\mathfrak a)\).
The same argument works for every derivative.

Choose \(t_0>0\) large enough that
\[
                 -t-\alpha<\ell,\qquad t+\alpha>d
                              \quad(t\ge t_0).
 \tag{MA-MO.561}
\]
The support of \(RU_th\) is
\([-t-\beta,-t-\alpha]\). For \(x>\ell\), every
argument \(x+\log N\mathfrak a\) is greater than that
support. Hence the literal inverse series proves
\[
             qF_{\bar\eta,RU_th}=0\qquad(t\ge t_0).
 \tag{MA-MO.562}
\]
This is exact vanishing of the dual source term before
Fourier transform, not an asymptotic estimate. The second
condition in (MA-MO.561) gives
\(\ell_{v_g}(U_th)=0\). Translation commutes with the
original inverse series, so
\(F_{\eta,U_th}(y)=F(y-t)\). Pairing the remaining
actual theta sum against \(v_g\) now gives
\[
 \begin{aligned}
 f(t)&=\ell_v(\mathscr D^{\,\rm raw}_{\eta,g}U_th)\\
     &=\int q(y)W(y)F(y-t)\,dy,\qquad t\ge t_0 .
 \end{aligned}
 \tag{MA-MO.563}
\]
For each such \(t\), the cutoff inverse is supported
between \(\ell\) and \(t+\beta\). Only finitely many
ideal shifts can meet the compact output test, so the
change of variables in this pairing is justified by
finite sums. It is the actual theta action of MA-CC.8.
The formula also shows that \(f\) is continuous on
every compact \(t\)-interval.

\emph{The actual paired Laplace transform.}
For \(\Re w>1/2\), define
\[
 \begin{aligned}
 \mathfrak F(w)&=\int_{t_0}^\infty e^{-wt}f(t)\,dt,\\
 B_v(w)&=\int \phi(y)W(y)e^{-wy}\,dy,\\
 \mathcal R_{h,v}(w)
    &=\int q(y)W(y)e^{-wy}
          \int_{y-t_0}^\infty e^{wu}F(u)\,du\,dy .
 \end{aligned}
 \tag{MA-MO.564}
\]
The first integral converges absolutely on that half-plane.
Indeed choose \(1/2<b<\Re w\) in (MA-MO.560). The
integrand in (MA-MO.563) is bounded by a constant times
\(e^{bt}\), and its \(y\)-interval has length at most
\(t+C\). This gives an integrable bound
\((t+C)e^{-(\Re w-b)t}\) for the double integral and
justifies Fubini.

Both \(B_v\) and \(\mathcal R_{h,v}\) are entire.
The function \(\phi W\) is compactly supported: \(W\)
is zero to the left of \(c\), and \(\phi\) is zero
sufficiently far right. For the second assertion put
\(y_0=\max(\ell,c)\). The nonzero part of the remainder
lies in
\[
       y_0\le y\le\beta+t_0,\qquad
                    y-t_0\le u\le\beta .
 \tag{MA-MO.565}
\]
Thus it is a compact double integral of a smooth
function times \(e^{w(u-y)}\), which proves entire
dependence with all derivatives. An empty interval
contributes zero.

Absolute Euler convergence on \(\Re w>1/2\) gives
\[
       \int F(u)e^{wu}\,du=\frac{H_h(w)}{A_\eta(w)},
       \qquad
       \int W(y)e^{-wy}\,dy=A_\eta(w)V_g(-w).
 \tag{MA-MO.566}
\]
The coefficients in the first equality are the literal
M\"obius coefficients; those in the second are the
literal forward theta coefficients. From
\[
 \int_{t_0}^\infty e^{-wt}F(y-t)\,dt
       =e^{-wy}\int_{-\infty}^{y-t_0}e^{wu}F(u)\,du
\]
and \(q=1-\phi\), equations (MA-MO.563)--(MA-MO.566)
therefore give the exact identity
\[
 \mathfrak F(w)
   =H_h(w)V_g(-w)
       -\frac{H_h(w)B_v(w)}{A_\eta(w)}
       -\mathcal R_{h,v}(w),
                    \qquad\Re w>1/2 .
 \tag{MA-MO.567}
\]
No narrow reciprocal multiplier is used to obtain it.

\emph{A cutoff coefficient survives every nonzero smoothing.}
At the fixed exterior point define
\[
 b_\lambda(x)=e^{-\lambda x}
       \sum_{\mathfrak a\ {\rm prime\ to}\ S}
         c_{\eta,\mathfrak a}(N\mathfrak a)^{-1/2-\lambda}
                                \phi(x+\log N\mathfrak a).
 \tag{MA-MO.568}
\]
The sum is locally finite. Let
\(r=\sup\{x:\phi(x)\ne0\}<\infty\). The function is
zero for \(x\ge r\). If \(r-\log2<x<r\), every
nonunit ideal has \(x+\log N\mathfrak a>r\); its
contribution is zero. Thus
\[
              b_\lambda(x)=e^{-\lambda x}\phi(x)
                              \quad(r-\log2<x<r).
 \tag{MA-MO.569}
\]
By the definition of \(r\), \(\phi\) is not identically
zero on that interval. Consequently \(b_\lambda\ne0\).

There is a direct exponential-order bound on its left
tail. Choose
\(\sigma_0>\max(1,1/2+\Re\lambda)\). For \(x\le r\),
only \(N\mathfrak a\le e^{r-x}\) contribute, so
\[
 \begin{aligned}
 |b_\lambda(x)|
 &\le C e^{-\Re\lambda x}
   e^{(r-x)(\sigma_0-1/2-\Re\lambda)}
                         \zeta_K(\sigma_0)\\
 &\le C_\lambda e^{-(\sigma_0-1/2)x}.
 \end{aligned}
 \tag{MA-MO.570}
\]
Its one-sided Mellin transform is therefore holomorphic
on \(\Re z>\sigma_0-1/2\). If \(C_gb_\lambda\) were
zero, absolute convolution on that half-plane would give
\[
                  \mathcal Mg(z)\,\mathcal Mb_\lambda(z)=0.
 \tag{MA-MO.571}
\]
The first factor is a nonzero entire function, by
Fourier uniqueness for the nonzero compact test \(g\).
The second would then vanish identically on that
half-plane by the identity theorem. On a fixed vertical
line there, Fourier uniqueness for the integrable function
\(e^{(\Re z)x}b_\lambda(x)\) would imply
\(b_\lambda=0\), a contradiction. Hence
\(C_gb_\lambda\ne0\), for arbitrary complex \(g\ne0\).

Choose \(v\in C_c^\infty\) with
\(\int v(x)C_gb_\lambda(x)\,dx\ne0\). Such a test
exists because the convolution is a nonzero smooth
function; one may multiply its complex conjugate by
a nonnegative compact bump in a region where it is
nonzero. Finite summation on the compact support of
\(v_g\) and the bilinear transpose identity give
\[
            B_v(\lambda)
               =\int v_g(x)b_\lambda(x)\,dx
               =\int v(x)C_gb_\lambda(x)\,dx\ne0 .
 \tag{MA-MO.572}
\]
This is why a possible zero of \(\mathcal Mg(\lambda)\)
does not remove the obstruction: the cutoff function
\(b_\lambda\), with its last support edge, is not a
pure exponential.

Finally choose a nonnegative nonzero
\(h_0\in C_c^\infty\) and \(h(x)=e^{-\lambda x}h_0(x)\).
Then \(H_h(\lambda)=\int h_0>0\). The preceding
identities hold for these \(h,v\), with \(t_0\)
increased to satisfy (MA-MO.561) if necessary.

\emph{A bounded translated orbit contradicts the actual pole.}
The elementary translation estimate in the exact sum norm is
\[
 \|U_th\|_H
    \le 2^{s/4}(1+t^2)^{s/4}e^{at}\|h\|_H
                                  \qquad(t\ge0).
 \tag{MA-MO.573}
\]
If (MA-MO.552) had a bounded \(H\)-extension, continuity
of \(\ell_v\) would imply
\(|f(t)|\le C e^{at}(1+t)^{s/2}\). Dominated
differentiation would then make \(\mathfrak F(w)\)
holomorphic throughout \(\Re w>a\). All other
nonreciprocal factors in (MA-MO.567) are entire.
Multiplying that identity by \(A_\eta(w)\) and using
the identity theorem on the connected half-plane
\(\Re w>a\) gives
\[
 A_\eta(w)\bigl[
      H_h(w)V_g(-w)-\mathcal R_{h,v}(w)-\mathfrak F(w)
            \bigr]=H_h(w)B_v(w)
                      \qquad(\Re w>a).
 \tag{MA-MO.574}
\]
At \(w=\lambda\), its left side is zero and its right
side is nonzero by construction. The finite Euler
deletions have no zeros on this half-plane, so a
completed zero of any multiplicity gives an actual
zero of \(A_\eta\) here. Equivalently the quotient
in (MA-MO.567) has a pole of that full multiplicity;
the numerator has been chosen nonzero. Thus no
higher-order cancellation has been assumed away.

The same argument shows more directly that
\[
       \sup_{t\ge t_0}
          \frac{|f(t)|}{e^{at}(1+t)^{s/2}}=\infty .
 \tag{MA-MO.575}
\]
Otherwise its finite bound alone would justify the same
holomorphic continuation and contradiction. Since \(f\)
is continuous on compact intervals, one can choose
\(t_n\to\infty\) for which that ratio tends to infinity.
Weighted Cauchy--Schwarz for \(\ell_v\) and
(MA-MO.573) imply (MA-MO.554).

\emph{The left side and the exact remaining boundary scope.}
The actual inverse Fourier pair and zero-singular
Poisson identity in (MA-MO.552) give, on the original
smooth sources,
\[
 \mathfrak S_\eta H_{\eta,y}=H_{\bar\eta,Ry},
 \qquad
 R\vartheta_\eta=\vartheta_{\bar\eta}\mathfrak S_\eta .
\]
Together with \(RC_g=C_{g^\vee}R\), these identities give
\[
           R\mathscr D^{\,\rm raw}_{\eta,g}R
                 =\mathscr D^{\,\rm raw}_{\bar\eta,g^\vee}.
 \tag{MA-MO.576}
\]
This uses both source sectors, not complex conjugation
of a vector. Reflection is an isometry of the sum
norm. A left-exterior completed zero of \(\eta\)
gives a right-exterior completed zero of \(\bar\eta\)
by the functional equation. Apply the proved case to
\(\bar\eta,g^\vee\) and then reflect the genuine inputs.
This proves (MA-MO.553)--(MA-MO.555).

Every part of the nonannihilation argument required
only \(g\ne0\). None required a nonzero value of its
transform at \(\lambda\), at either polar exponent,
or a nonzero \(g(0)\). The claimed test subclasses
therefore follow.
\end{proof}

\paragraph{An every-width operator criterion, with its arithmetic input exposed.}
For this fixed nontrivial character the following statements
are equivalent:
\[
 \begin{array}{ll}
 {\rm (i)}&
   \Re\rho=1/2\quad\text{for every nontrivial completed zero};\\
 {\rm (ii)}&
   \forall\,0<a<1/2,\ \exists s\ge0,\
                          \exists\,0\ne g\in C_c^\infty:\\
 &\quad\mathscr D^{\,\rm raw}_{\eta,g}
       \text{ has a bounded }H_{\Sigma,a,s}\text{ extension};\\
 {\rm (iii)}&
   \forall\,0<a<1/2,\ \forall s\ge0,\
                              \forall g\in C_c^\infty:\\
 &\quad\mathscr D^{\,\rm raw}_{\eta,g}
       \text{ has an }\mathcal S_1(H_{\Sigma,a,s})\text{ extension}.
 \end{array}
 \tag{MA-MO.577}
\]
The first statement and conjugation give the premise of
MA-MO-022 for both characters with \(a_*=0\), proving
the third. The third implies the second. The second
and (MA-MO.555), for arbitrarily small \(a\), force
the first. This is an actual original-operator criterion;
it asserts none of these arithmetic antecedents anew.
The second statement can also restrict its nonzero tests
by \(g(0)=\mathcal Mg(\pm1/2)=0\). That class is
nonempty: for nonzero \(h_0\in C_c^\infty((1,2))\),
\(g=(\partial_x^2-1/4)h_0\) has these properties and
cannot be zero by the elementary constant-coefficient ODE.

At a fixed width, containment of the entire divisor in
a strictly smaller closed strip is sufficient by
MA-MO-022, and (MA-MO.555) is necessary.
The present Laplace argument is holomorphic only on
\(\Re w>a\). It does not settle zeros on the boundary,
their polynomial-weight thresholds, or sufficiency when
zeros approach an edge without a uniform margin.

\paragraph{Original-source and framework placement.}
The only source identities used are the literal ideal
Euler sums and inverse coefficients (MA-CC.1 and .8),
the structural coinvariant inclusion (MA-MO.470,.476),
the original \(p_+\) formula, and nontrivial-sector
Poisson summation. Their primary locations are Meyer
(2005), supplied \texttt{main.tex} lines 3592--3707,
2630--2675, and 3808--4043. The proof uses the
unconditional source identities from MA-MO-022; it does
not import that theorem's conditional narrow inverse
identification. The finite ideal sums and the cleared
Laplace identity above provide the new comparison.

No existence of a strictly exterior arithmetic zero is
asserted. The theorem proves what such a zero would do
to these actual source maps, rather than only to an
uncontinued scalar trace or a bare reciprocal multiplier.
The source restriction and quotient-return maps in
MA-MO-012/017 and MA-PR remain different domains.
In the inherited cote 1 framework, scans 079--080 and
111--112 through BR-TR-01/02, an exterior divisor factor
here prevents a bounded realization of this particular
full source action. Its finite-length algebraic class
does not repair the missing analytic norm estimate.
No positivity, additional motivic weight, or independent
FT002 deficiency vanishing follows from (MA-MO.577).
